% Resumen en inglés
\chapter*{Abstract}

\begin{abstractEn}
The increasing frecuency of cancer cases, as well as its severity, has made the cientific comunity amplify its efforts on researching this illness. Due to this increased interest, numerous tools have appeared to aid with this purpose, some of them, focused on predicting the evolution of cancerous tumours, are known as tumoral progression simulators.

On this project, the functionality of one of these simulators, \textbf{OncoSimulR} is going to be extended. A group of changes is going to be added to the program in order to allow the user to specify arbitrarily complex user variables that depend on existing parametesrs of the simulation, giving the user the possibility of tracking any data of the tumoral progression. This variables will also be usable for the definition of the interventions that the program can already emulate, adapting them to the current state of the tumor in a given moment, in order to achieve the simulation of adaptive therapy.

Furthermore, some tests will be developed in order to ensure the optimal performance of this new functionality. Lastly, examples of some use cases of the program will be given, illustrating the new possible ways to use the tool.
\end{abstractEn}

% Palabras clave en inglés
\begin{keywordsEn}
Cancer, tumor, tumoral progression simulator, OncoSimulR, user variables, interventions, adaptive therapy.
\end{keywordsEn}

% Resumen en español
\chapter*{Resumen}

\begin{abstractEs}
La creciente preocupación por la proliferación del \gls{cancer} como enfermedad, así como por la gravedad de este, ha propiciado que en los últimos años se hayan multiplicado los esfuerzos en su investigación. Con este fin han surgido numerosas herramientas para predecir la evolución de un \gls{cancer}, conocidas cómo \glspl{simulador}.

En este trabajo se va a ampliar la funcionalidad de uno de estos simuladores, \textbf{OncoSimulR}. Sobre éste se va a introducir la posibilidad de que el usuario defina una serie de variables de usuario tan complejas como desee, en función de parámetros de la simulación, que permitirán al usuario obtener prácticamente cualquier dato que quiera acerca de la evolución del \gls{tumor}. Además se va a permitir que estas variables se empleen en la definición de las intervenciones que esta aplicación permite reproducir, con el fin de adaptar estas al estado del \gls{tumor} en cada momento, logrando de este modo emular la conocida como terapia adaptativa.

Tras explicar en qué consisten los cambios llevados a cabo para esto, se describirá la batería de pruebas empleada para comprobar su correcto funcionamiento, y se mostrarán algunos ejemplos prácticos del uso del programa, ilustrando esta nueva funcionalidad.
\end{abstractEs}

% Palabras clave en español
\begin{keywordsEs}
\Gls{cancer}, \gls{tumor}, \gls{simulador}, OncoSimulR, variables de usuario, intervenciones,terapia adaptativa.
\end{keywordsEs}
